%! Author = Saeid Yazdani
%! Date = 1/8/2024


\section{Will Built-In Test Render \acp{ATE} Obsolete?}
\label{sec:bist-vs-ate}

In the \autoref{subsec:dfx-dft}, some of the significant testing trends
in the industry were presented.
It was established that built-in test facilities such as scan chains and
\ac{BIST} are commonly used methods for reducing test time and effort.
On the other hand, built-in structures including derivatives and similar
concepts, by no means are supposed or made it possible \emph{yet} to
remove \acp{ATE} from
the \emph{semiconductor ecosystem} (\autoref{fig:semiconductor-ecosystem}).

\acp{ATE} are still mandatory for powering up and applying test sequences
(either directly or by activating scan chains or \ac{BIST} cores
on the \ac{DUT}) and log and verify the results.

In addition, \ac{ATE} based tests are still carried out at wafer
and package levels to obtain pass/fail results (i.e. whether or not the chip is
functional), and performance-class binning of the chips for different maximum
clock speeds, which is a common practice in
consumer-grade products like \ac{CPU} and  \ac{GPU} chips.

\subsection{Parametric tests and Influence on ATE-DUT Power Supplies}

Although \ac{BIST} (\autoref{subsec:bist}) has improved
over time,
performing parametric tests is an area where traditional \acp{ATE} are still
dominating as most of the tests are inherently analog and
require precision measurement capabilities. Such tests include:

\begin{itemize}[topsep=8pt,itemsep=4pt,partopsep=4pt, parsep=4pt]

    \item \textbf{\emph{Voltage and Current Measurements}:}
    Stimulating and measuring the response of the \ac{DUT} for known effects
    such as leakage currents, saturation voltages, driver capabilities, and
    voltage thresholds.

    \item \textbf{\emph{Power Consumption}:}
    Bring the \ac{DUT} into known states and measure power
    consumption (static and dynamic power), and check if the
    power consumption is within predefined limits.

    \item \textbf{\emph{Environmental tests}:}
    \acp{ATE} with thermal modules bring
    \ac{DUT} to high or low temperatures and perform tests for simulating
    harsh application environments.

    \item \textbf{\emph{Analog parameters}:}
    \acp{ATE} can measure gain, bandwidth, \ac{SNR}, harmonic distortions,
    and etc.

    \item \textbf{\emph{\ac{RF} Evaluation}:}
    For \ac{RF} mixed-signal and \ac{SoC} devices and transceivers,
    measurements of parameters such as frequency response and
    modulation characteristics are vital.
\end{itemize}

Many parametric tests can be performed relatively quickly and with standard accuracy. 
However, measurements involving high-resolution \ac{ADC}s/\ac{DAC}s or very small currents 
and voltages often require specialized precision equipment and longer test times, which can 
limit the cost-effectiveness of conventional \acp{ATE}.